% 数学ノート集・共通テンプレート v2（積分論2026年度第1部の版面を参考）
% A5判。章内で定義・定理・例・注意等の番号を共通化。
% コンパイル：platex main.tex（3回），dvipdfmx main.dvi
%              dvipdfmx math_note_template_v2.dvi
\documentclass[a5paper,10pt,onecolumn]{jsarticle}
\usepackage[dvipdfmx,top=22mm,bottom=20mm,inner=19mm,outer=18mm,headheight=13pt,headsep=6mm]{geometry}
\usepackage{amsmath,amssymb,amsthm,mathtools}
\usepackage{enumitem,fancyhdr,titlesec}
\usepackage[dvipdfmx,bookmarks=true,colorlinks=false]{hyperref}
\usepackage{pxjahyper}

% ----- 差し替える文書情報 -----
\newcommand{\NoteTitle}{開写像定理・有界逆写像定理・閉グラフ定理}
\newcommand{\NoteSeries}{数学ノート集}

% ----- 文字・組版 -----
% 参考PDFの質素な明朝＋Computer Modernの印象を引き継ぎながら，
% A5判で窮屈にならないよう本文の行間を少し確保する．
\linespread{1.13}
\setlength{\parindent}{1zw}
\setlength{\parskip}{0pt}
\setlength{\abovedisplayskip}{6pt plus 2pt minus 2pt}
\setlength{\belowdisplayskip}{6pt plus 2pt minus 2pt}
\setlength{\abovedisplayshortskip}{3pt plus 2pt minus 1pt}
\setlength{\belowdisplayshortskip}{4pt plus 2pt minus 1pt}
\setlist[enumerate]{topsep=3pt,itemsep=1pt,parsep=0pt,leftmargin=2.4zw}
\titleformat{\section}{\normalfont\normalsize\bfseries}{\thesection}{0.8zw}{}
\titlespacing*{\section}{0pt}{2.0ex plus .4ex}{.9ex}
\titleformat{\subsection}{\normalfont\normalsize\bfseries}{\thesubsection}{0.7zw}{}
\titlespacing*{\subsection}{0pt}{1.5ex plus .3ex}{.6ex}
\raggedbottom

% ----- 柱とページ番号：ページ上部、罫線なし -----
\fancypagestyle{mathnote}{%
  \fancyhf{}%
  \fancyhead[L]{\footnotesize \NoteSeries\ （\NoteTitle）}%
  \fancyhead[R]{\footnotesize \thepage}%
  \renewcommand{\headrulewidth}{0pt}%
  \renewcommand{\footrulewidth}{0pt}%
}
\pagestyle{mathnote}

% ----- 節内通し番号の定義・定理群 -----
\newtheoremstyle{mathnotetheorem}
 {4.5pt}{4.5pt}{\normalfont}{}%
 {\bfseries}{}{0.65zw}%
 {\thmname{#1}\thmnumber{ #2}\thmnote{（#3）}}
\theoremstyle{mathnotetheorem}
\newtheorem{theorem}{定理}[section]
\newtheorem{definition}[theorem]{定義}
\newtheorem{lemma}[theorem]{補題}
\newtheorem{proposition}[theorem]{命題}
\newtheorem{corollary}[theorem]{系}
\newtheorem{example}[theorem]{例}
\newtheorem{remark}[theorem]{注意}
\renewcommand{\proofname}{証明}
\renewcommand{\qedsymbol}{$\square$}
\makeatletter
\renewenvironment{proof}[1][\proofname]{%
  \par\pushQED{\qed}\normalfont\topsep4\p@\@plus2\p@\relax
  \trivlist\item[\hskip\labelsep\bfseries #1]\ignorespaces
}{\popQED\endtrivlist\@endpefalse}
\makeatother

% ----- 使うノートで適宜拡張する記号 -----
\newcommand{\R}{\mathbb{R}}
\newcommand{\norm}[1]{\lVert #1\rVert}
\newcommand{\abs}[1]{\lvert #1\rvert}

\begin{document}
\section*{読了後にできること}
本稿を読むための前提は，開写像定理，有界逆写像定理，Banach空間．
本稿の具体的な到達課題は，完備性が必要なことを具体的な反例で確認する．

\section{Banach空間とBaireのカテゴリー定理}

以下，$\mathbb K$ は実数体または複素数体とする．
ノルム空間 $X$ の開球を
$B_X(x,r)=\{u\in X:\|u-x\|<r\}$ と書く．
$X$ がノルムによる距離に関して完備なら
\emph{Banach空間}という．
以下，ノルム空間の間の有界線形写像全体を $\mathcal B(X,Y)$ と書く．
一般の距離空間 $(Z,d)$ における開球は
$B_Z(z,r)=\{w\in Z:d(z,w)<r\}$ と表し，
空間が明らかな場合には添字を省略する．

\begin{definition}
連続線形写像 $T:X\to Y$ が\emph{開写像}であるとは，
$X$ の任意の開集合 $U$ に対して
$T(U)$ が $Y$ で開であることをいう．
\end{definition}

\begin{lemma}[Baireのカテゴリー定理]\label{lem:baire}
空でない完備距離空間 $Z$ の可算個の稠密開集合
$U_1,U_2,\ldots$ に対し，
$\bigcap_{n\ge1}U_n$ は稠密である．
\end{lemma}
\begin{proof}
任意の空でない開集合 $V\subset Z$ を取る．
$U_1$ が稠密開なので，
ある $z_1\in V\cap U_1$ と $0<r_1<2^{-1}$ が存在して
$\overline{B(z_1,r_1)}\subset V\cap U_1$ とできる．
同様に，$z_n,r_n$ が得られているとき，
$U_{n+1}\cap B(z_n,r_n)$ が空でない開集合であることから，
\[
 \overline{B(z_{n+1},r_{n+1})}
 \subset U_{n+1}\cap B(z_n,r_n),\qquad
 0<r_{n+1}<2^{-(n+1)}
\]
となるように選べる．
$m>n$ なら $d(z_m,z_n)<r_n<2^{-n}$ であるため
$(z_n)$ はCauchy列となる．完備性から $z_n\to z$ を得る．
固定した $n$ について $m>n$ の $z_m$ は
$\overline{B(z_n,r_n)}$ に属するので，
その閉性により
$z\in\overline{B(z_n,r_n)}\subset U_n$ となる．
また $z\in V$ である．従って
$V\cap\bigcap_{n\ge1}U_n\ne\varnothing$ となり，稠密である．
\end{proof}

\begin{corollary}[閉集合被覆の補題]\label{cor:cover}
空でない完備距離空間 $Z$ を
閉集合 $F_1,F_2,\ldots$ で被覆するとき，
少なくとも一つの $F_n$ は空でない内部を持つ．
\end{corollary}
\begin{proof}
全ての $F_n$ の内部が空なら，
$Z\setminus F_n$ は稠密開である．
補題\ref{lem:baire}によって
$\bigcap_n(Z\setminus F_n)$ は稠密であり，特に空でない．
これは $\bigcup_nF_n=Z$ と矛盾する．
\end{proof}

\begin{lemma}[スケーリングと開写像性]\label{lem:openball}
$T:X\to Y$ を線形写像とする．
ある $r>0$ と $\delta>0$ に対して
\[
 B_Y(0,\delta)\subset T(B_X(0,r))
\]
が成立すれば $T$ は開写像である．
\end{lemma}
\begin{proof}
任意の開集合 $U\subset X$ と $y_0=Tx_0\in T(U)$ を取る．
ある $\varepsilon>0$ が存在して
$x_0+B_X(0,\varepsilon)\subset U$ となる．
線形性とスケーリングから
\[
 T(B_X(0,\varepsilon))
 \supset B_Y(0,\delta\varepsilon/r).
\]
従って
$y_0+B_Y(0,\delta\varepsilon/r)\subset T(U)$ となり，
$T(U)$ は開である．
\end{proof}

\section{開写像定理の証明}

\begin{lemma}[閉包の球から真の球へ]\label{lem:closureball}
$X$ をBanach空間，$Y$ をノルム空間，
$T\in\mathcal B(X,Y)$ とする．
ある $\delta>0$ に対して
\[
 B_Y(0,\delta)
 \subset\overline{T(B_X(0,1))}
 \tag{$*$}
\]
が成立するなら，
\[
 B_Y(0,\delta/2)\subset T(B_X(0,1))
\]
が成り立つ．
\end{lemma}
\begin{proof}
$y\in B_Y(0,\delta/2)$ を固定する．
$(*)$ を正数 $a$ 倍すると，
\[
 B_Y(0,a\delta)\subset
 \overline{T(B_X(0,a))}\qquad(a>0)
 \tag{$**$}
\]
が従う．

まず $(**)$ を $a=1/2$ に適用し，
$\|x_1\|<2^{-1}$ かつ
$\|y-Tx_1\|<\delta 2^{-2}$ を満たす
$x_1\in X$ を取る．
帰納的に
\[
 \|x_j\|<2^{-j}\quad(1\le j\le n),
 \qquad
 \left\|y-T\sum_{j=1}^{n}x_j\right\|
 <\delta 2^{-(n+1)}
\]
が成立すると仮定する．
残差 $r_n=y-T\sum_{j=1}^nx_j$ は
$B_Y(0,\delta2^{-(n+1)})$ に属する．
$(**)$ を $a=2^{-(n+1)}$ に適用すれば，
$\|x_{n+1}\|<2^{-(n+1)}$ かつ
\[
 \|r_n-Tx_{n+1}\|<\delta2^{-(n+2)}
\]
となる $x_{n+1}$ が取れる．
このように $(x_n)$ を構成する．
$\sum_{n=1}^\infty\|x_n\|\le\sum_{n=1}^\infty2^{-n}=1$
だから，Banach空間 $X$ の完備性により
$x=\sum_{n=1}^\infty x_n\in X$ が存在して $\|x\|\le1$ となる．
また，$T$ の連続性と残差の評価から
\[
 Tx=\lim_{n\to\infty}T\sum_{j=1}^nx_j
 =\lim_{n\to\infty}(y-r_n)=y
\]
である．

最後に $\|y\|<\delta/2$ なので，
$(1+\eta)\|y\|<\delta/2$ を満たす $\eta>0$ を取れる．
上の構成を $y'=(1+\eta)y$ に適用すると，
$Tx'=y'$，$\|x'\|\le1$ を満たす $x'\in X$ が得られる．
従って $z=x'/(1+\eta)$ は $\|z\|<1$ および
$Tz=y$ を満たす．これで結論を得る．
\end{proof}

\begin{theorem}[開写像定理]\label{thm:open}
$X,Y$ をBanach空間とし，
$T:X\to Y$ を全射な有界線形作用素とする．
このとき $T$ は開写像である．
\end{theorem}
\begin{proof}
$B=B_X(0,1)$ と置く．$T$ は全射なので
\[
 Y=\bigcup_{n=1}^\infty T(nB)
 \subset\bigcup_{n=1}^\infty\overline{T(nB)}
 \subset Y,
\]
従って
$Y=\bigcup_{n=1}^\infty\overline{T(nB)}$ である．
各 $\overline{T(nB)}$ は閉集合であり，
$Y$ は完備であるから，
系\ref{cor:cover}によって
ある正整数 $N$ と $y_0\in Y$, $r>0$ が存在して
\[
 B_Y(y_0,r)\subset\overline{T(NB)}
\]
となる．特に $y_0\in\overline{T(NB)}$ である．

任意の $v\in B_Y(0,r)$ について
$y_0+v$ と $y_0$ はともに
$\overline{T(NB)}$ に属するから，
\[
 v=(y_0+v)-y_0
 \in\overline{T(NB)}-\overline{T(NB)}
 \subset\overline{T(2NB)}.
\]
最後の包含は $NB-NB\subset2NB$ と，
閉包における減法の連続性から従う．
スケーリングによって
\[
 B_Y(0,r/(2N))
 \subset\overline{T(B)}.
\]
補題\ref{lem:closureball}を
$\delta=r/(2N)$ に適用すると
\[
 B_Y(0,r/(4N))\subset T(B).
\]
補題\ref{lem:openball}によって $T$ は開写像である．
\end{proof}

\begin{remark}
開写像定理には定義域と値域の両方のBanach空間としての
完備性が用いられる．
証明で $Y$ の完備性は閉集合被覆の補題に，
$X$ の完備性は近似解の級数の収束に使われた．
\end{remark}

\section{有界逆写像定理}

\begin{theorem}[有界逆写像定理]\label{thm:inverse}
$X,Y$ をBanach空間とし，
$T\in\mathcal B(X,Y)$ を全単射とする．
このとき逆写像 $T^{-1}:Y\to X$ は有界線形作用素である．
\end{theorem}
\begin{proof}
$T$ が線形な全単射なので $T^{-1}$ は線形である．
開写像定理によって $T(B_X(0,1))$ は
$Y$ の $0$ を含む開集合である．
従ってある $\delta>0$ に対し
$B_Y(0,\delta)\subset T(B_X(0,1))$ となる．
$y\ne0$ に対して
$v=\delta y/(2\|y\|)$ と置けば
$\|v\|<\delta$ だから
$\|T^{-1}v\|<1$ である．線形性から
\[
 \|T^{-1}y\|
 <\frac{2}{\delta}\|y\|.
\]
$y=0$ は自明なので $T^{-1}$ は有界である．
\end{proof}

\begin{corollary}[Banach空間のノルムの同値性]
線形空間 $V$ に二つのノルム $\|\cdot\|_1,\|\cdot\|_2$
が与えられ，両方がBanach空間を定めるとする．
ある $C>0$ について
$\|x\|_2\le C\|x\|_1$ がすべての $x\in V$ で
成立するなら，ある $D>0$ に対し
\[
 \|x\|_1\le D\|x\|_2\qquad(x\in V)
\]
が成立する．
\end{corollary}
\begin{proof}
恒等写像
$I:(V,\|\cdot\|_1)\to(V,\|\cdot\|_2)$ は
有界線形全単射である．定理\ref{thm:inverse}によって
$I^{-1}$ が有界となるので結論を得る．
\end{proof}

\section{商空間と全射作用素}

\begin{definition}
ノルム空間 $X$ の閉線形部分空間 $M$ に対して
線形商空間 $X/M$ に
\[
 \|x+M\|_{X/M}=\inf_{m\in M}\|x-m\|_X
\]
を定め，これを\emph{商ノルム}という．
自然な商写像を $\pi:X\to X/M$ と書く．
\end{definition}

\begin{lemma}\label{lem:quotient-norm}
上の式は $X/M$ にノルムを定め，
$\|\pi x\|_{X/M}\le\|x\|_X$ が成立する．
\end{lemma}
\begin{proof}
代表元を $x$ から $x+m_0$ に替えても，
$\inf_{m\in M}\|x+m_0-m\|$ は同じであるから
well-definedである．
斉次性は $a\ne0$ に対して
\[
 \inf_{m\in M}\|ax-m\|
 =|a|\inf_{m'\in M}\|x-m'\|
\]
から従い，$a=0$ は自明である．
三角不等式は
$\|x+y-(m+n)\|\le\|x-m\|+\|y-n\|$
で $m,n\in M$ に関する下限を取ればよい．
$\|x+M\|=0$ なら $x\in\overline M=M$ であるから，
$X/M$ における零元となる．
最後の不等式は $m=0$ を取ることで従う．
\end{proof}

\begin{proposition}[Banach商空間の完備性]\label{prop:quotient}
$X$ をBanach空間，$M$ を閉線形部分空間とすると，
商ノルムを入れた $X/M$ はBanach空間である．
\end{proposition}
\begin{proof}
$(\xi_n)$ を $X/M$ のCauchy列とする．
部分列 $(\xi_{n_k})$ を
\[
 \|\xi_{n_{k+1}}-\xi_{n_k}\|_{X/M}<2^{-k}
 \qquad(k\ge1)
\]
となるように選ぶ．
各差の代表元 $u_k\in X$ を
$\pi u_k=\xi_{n_{k+1}}-\xi_{n_k}$，
$\|u_k\|_X<2^{-k+1}$ となるように取ることができる．
実際，商ノルムの下限の定義により，
与えられた代表元を $M$ の元だけ調整すればよい．
$\xi_{n_1}$ の代表元 $x_1\in X$ を選び，
\[
 x=x_1+\sum_{k=1}^\infty u_k
\]
と置く．級数はノルム絶対収束し，
$X$ の完備性により $x\in X$ が存在する．
自然な商写像 $\pi$ は連続なので
\[
 \pi x=\lim_{j\to\infty}
 \left(\xi_{n_1}+\sum_{k=1}^{j-1}
 (\xi_{n_{k+1}}-\xi_{n_k})\right)
 =\lim_{j\to\infty}\xi_{n_j}.
\]
従って部分列は $\pi x$ に収束する．
Cauchy列で収束部分列を持つものは全体が同じ極限に収束するので，
$\xi_n\to\pi x$ となり完備である．
\end{proof}

\begin{theorem}[全射作用素と商空間]\label{thm:quotient-isom}
$X,Y$ をBanach空間，$T:X\to Y$ を全射有界線形作用素とし，
$M=\ker T$ と置く．このとき
\[
 \widetilde T:X/M\longrightarrow Y,\qquad
 \widetilde T(x+M)=Tx
\]
は有界線形全単射であり，逆写像も有界である．
特に，ある $C>0$ が存在して
\[
 \inf_{m\in\ker T}\|x-m\|
 \le C\|Tx\|\qquad(x\in X)
\]
が成立する．
\end{theorem}
\begin{proof}
$M$ は連続写像の核であるから閉であり，
命題\ref{prop:quotient}によって $X/M$ はBanach空間である．
$Tx=Tx'$ と $x-x'\in M$ は同値なので
$\widetilde T$ はwell-definedで単射である．
$T$ の全射性から $\widetilde T$ は全射である．
任意の $m\in M$ について
\[
 \|\widetilde T(x+M)\|
 =\|T(x-m)\|\le\|T\|\,\|x-m\|.
\]
下限を取って
$\|\widetilde T(x+M)\|\le\|T\|\|x+M\|$ を得る．
定理\ref{thm:inverse}から
$\widetilde T^{-1}$ は有界であり，
\[
 \|x+M\|
 =\|\widetilde T^{-1}(Tx)\|
 \le\|\widetilde T^{-1}\|\,\|Tx\|
\]
である．
\end{proof}

\begin{remark}
全射 $T$ に対して，$Tx=y$ を満たす $x$ は一般に一意ではない．
上の評価は，解の集合が核の剰余類であることを踏まえたものである．
この式は $T$ の開写像性を定量的に表している．
\end{remark}

\section{閉グラフ定理}

\begin{definition}
線形作用素 $T:D(T)\to Y$ の\emph{グラフ}を
\[
 G(T)=\{(x,Tx):x\in D(T)\}\subset X\times Y
\]
と定める．ここで $D(T)$ は $X$ の線形部分空間とする．
$G(T)$ が $X\times Y$ の閉部分集合なら
$T$ を\emph{閉作用素}という．
\end{definition}

\begin{lemma}[直積の完備性]\label{lem:product}
$X,Y$ をBanach空間とし，
$X\times Y$ に $\|(x,y)\|=\|x\|_X+\|y\|_Y$
を定めるとBanach空間となる．
また，Banach空間の閉線形部分空間はBanach空間である．
\end{lemma}
\begin{proof}
直積のCauchy列 $((x_n,y_n))$ については，
各成分がそれぞれCauchy列である．
完備性から $x_n\to x$，$y_n\to y$ が成り立ち，
定義した直積ノルムに関して
$(x_n,y_n)\to(x,y)$ となる．
閉線形部分空間 $M$ のCauchy列は元のBanach空間でもCauchyである．
元の空間での極限は $M$ の閉性によって $M$ に属する．
\end{proof}

\begin{theorem}[閉グラフ定理]\label{thm:closedgraph}
$X,Y$ をBanach空間とし，
$T:X\to Y$ を全域で定義された線形作用素とする．
このとき次は同値である．
\begin{enumerate}
\item $T$ は有界である．
\item $G(T)$ は $X\times Y$ の閉集合である．
\end{enumerate}
\end{theorem}
\begin{proof}
$T$ が有界なら連続である．
$(x_n,Tx_n)\to(x,y)$ とすると
$x_n\to x$ および $Tx_n\to y$ である．
連続性によって $Tx_n\to Tx$ でもあるから
$y=Tx$ となり，$(x,y)\in G(T)$ である．
距離空間では点列による閉性の特徴づけが使えるため，
$G(T)$ は閉である．

逆に $G(T)$ が閉とする．
$G(T)$ は線形部分空間であるから，
補題\ref{lem:product}により
直積ノルムでBanach空間となる．
第一射影
\[
 P:G(T)\to X,\qquad P(x,Tx)=x
\]
は連続線形作用素であり，
$\|P(x,Tx)\|\le\|(x,Tx)\|$ となる．
$T$ が $X$ の全域に定義されているから $P$ は全射，
グラフの定義から単射である．
定理\ref{thm:inverse}より
逆写像
$P^{-1}:X\to G(T)$ は有界である．
第二射影 $Q:G(T)\to Y$，
$Q(x,Tx)=Tx$ も有界である．
従って
$T=Q\circ P^{-1}$ は有界である．
\end{proof}

\begin{corollary}[閉グラフの判定法]
$X,Y$ がBanach空間で，
$T:X\to Y$ が全域で定義された線形作用素なら，
次の条件は $T$ の有界性と同値である：
\[
 x_n\to x\text{ in }X,\quad
 Tx_n\to y\text{ in }Y
 \quad\Longrightarrow\quad Tx=y.
\]
\end{corollary}
\begin{proof}
条件は $G(T)$ が点列について閉であることと同値である．
$X\times Y$ は距離空間なので，
点列について閉であることと閉集合であることは同値である．
定理\ref{thm:closedgraph}を適用する．
\end{proof}

\section{閉だが非有界な微分作用素}

\begin{example}[微分作用素]\label{ex:derivative}
$X=Y=C([0,1],\mathbb K)$ に一様ノルムを入れる．
\[
 D(T)=C^1([0,1],\mathbb K),\qquad Tf=f'
\]
と定める．
このとき $T:D(T)\to Y$ は閉作用素であるが，
$D(T)$ に $X$ からの一様ノルムを入れたとき有界でない．
\end{example}
\begin{proof}
まず閉性を示す．
$(f_n,f_n')\to(f,g)$ が
$X\times Y$ の直積ノルムで成り立つとする．
すなわち $f_n\to f$ と $f_n'\to g$ はともに一様収束する．
微積分学の基本定理により
\[
 f_n(t)-f_n(0)=\int_0^t f_n'(s)\,ds
\]
である．両辺で $n\to\infty$ とすると
一様収束により極限と積分を交換でき，
\[
 f(t)-f(0)=\int_0^t g(s)\,ds
\]
を得る．$g$ は連続なので右辺は $C^1$ 関数であり，
$f\in C^1([0,1])$ かつ $f'=g$ となる．
従って $(f,g)\in G(T)$ であり，$G(T)$ は閉である．

一方 $f_n(t)=t^n$ と置けば
$\|f_n\|_\infty=1$ であるのに
$\|Tf_n\|_\infty=n$ である．
従って $T$ は有界でない．
\end{proof}

\begin{remark}
この例は閉グラフ定理に矛盾しない．
$T$ はBanach空間 $X$ の全域に定義されておらず，
$D(T)=C^1([0,1])$ は一様ノルムについて完備でないからである．
たとえば $|t-\tfrac12|$ は
$C^1$ 関数で一様近似できるが自身は $C^1$ でない．
\end{remark}

\begin{proposition}[グラフノルム]\label{prop:graphnorm}
$X,Y$ をBanach空間，$T:D(T)\subset X\to Y$
を線形作用素とする．
\[
 \|x\|_T=\|x\|_X+\|Tx\|_Y\qquad(x\in D(T))
\]
と置く．このとき $T$ が閉作用素であることと
$(D(T),\|\cdot\|_T)$ がBanach空間であることは同値である．
\end{proposition}
\begin{proof}
写像
$J:D(T)\to X\times Y$，
$Jx=(x,Tx)$ は線形であり，
$\|Jx\|_{X\times Y}=\|x\|_T$ である．
従って $J$ は $(D(T),\|\cdot\|_T)$ を
直積の部分空間 $G(T)$ に等距離に写す全単射である．
$X\times Y$ はBanach空間だから，
$G(T)$ が閉であることと
$G(T)$ が部分空間のノルムで完備であることが同値となる．
実際，閉ならば閉部分空間の完備性によって完備であり，
逆に完備部分空間内で収束する点列はCauchyなので，
その極限が再び部分空間に属することから閉となる．
\end{proof}

\section{完備性を外した場合}

\begin{example}[開写像性と逆写像の有界性の失敗]
線形空間 $V=C^1([0,1],\mathbb K)$ に
\[
 \|f\|_1=\|f\|_\infty+\|f'\|_\infty,\qquad
 \|f\|_0=\|f\|_\infty
\]
の二つのノルムを入れる．
恒等写像
\[
 I:(V,\|\cdot\|_1)\longrightarrow(V,\|\cdot\|_0)
\]
は連続線形全単射だが，逆写像は有界でない．
従って $I$ は開写像でない．
\end{example}
\begin{proof}
$\|f\|_0\le\|f\|_1$ だから $I$ は有界であり，
恒等写像なので全単射である．
一方 $f_n(t)=t^n$ に対して
$\|f_n\|_0=1$，$\|f_n\|_1=1+n$ なので
$I^{-1}$ は有界でない．
$I$ が開写像なら $I^{-1}$ は連続となるから，
$I$ は開写像ではない．
\end{proof}

\begin{proposition}
上の例で $(V,\|\cdot\|_1)$ はBanach空間だが，
$(V,\|\cdot\|_0)$ はBanach空間でない．
\end{proposition}
\begin{proof}
$(f_n)$ を $\|\cdot\|_1$ に関するCauchy列とすると，
$(f_n)$ と $(f_n')$ はそれぞれ一様ノルムについて
Cauchy列である．$C([0,1])$ の完備性により
$f_n\to f$，$f_n'\to g$ が一様に成り立つ．
前の微分作用素の閉性の証明と同様に
$f\in C^1$，$f'=g$ となり，
$\|f_n-f\|_1\to0$ が従う．

他方，連続だが微分不可能な関数
$h(t)=|t-\tfrac12|$ は，
例えば
\[
 h_n(t)=\sqrt{(t-\tfrac12)^2+n^{-2}}
\]
によって一様近似される．実際
$0\le h_n(t)-h(t)\le1/n$ であり，
各 $h_n$ は $C^1$ である．
従って $(h_n)$ は $\|\cdot\|_0$ に関する
Cauchy列であるが，$V$ 内に極限を持たない．
\end{proof}

\begin{remark}
定理\ref{thm:open}と定理\ref{thm:inverse}は
Banach空間における完備性を不可欠な仮定としている．
定理\ref{thm:closedgraph}で線形作用素が全域で
定義されるという仮定も同様に重要である．
\end{remark}


\section{到達課題の完全解答}
\noindent\textbf{逆写像定理で完備性が必要な例．}
\begin{proposition}
$C^1([0,1])$ にノルム
$\|f\|_{C^1}=\|f\|_\infty+\|f'\|_\infty$ と
$\|f\|_\infty$ をそれぞれ入れる．
恒等写像
\[
I:(C^1,\|\cdot\|_{C^1})
\longrightarrow(C^1,\|\cdot\|_\infty)
\]
は有界線形全単射だが，
逆写像は有界でない．
\end{proposition}
\begin{proof}
$\|If\|_\infty\le\|f\|_{C^1}$ だから有界であり，
集合として恒等写像なので全単射である．
$f_n(x)=\sin(nx)/n$ に対し
$\|f_n\|_\infty\le1/n\to0$ だが，
$\|f_n'\|_\infty=1$ なので
$\|f_n\|_{C^1}\ge1$．
従って逆写像は連続でない．
値域の $\|\cdot\|_\infty$ では $C^1$ は
完備でないため，有界逆写像定理の仮定に反しない．
\end{proof}
\noindent\textbf{到達確認．}
$\|f\|_{C^1}$ による $C^1([0,1])$ はBanach空間か．
\par\noindent\textbf{解答．}
$C^1$-Cauchy列では $f_n,f_n'$ が
それぞれ一様Cauchyとなる．一様極限を $f,g$ とすると
積分公式から $f'=g$ が従うので完備である．
\end{document}
